\section{Проверка работоспособности системы МОТ в различных условиях}

Для проверки работоспособности разработанной системы мультиобъектного отслеживания необходимо провести экспериментов. В качестве экспериментов будут выступать видео с различными условиями. В качестве условий были выделены следующие критерии:

\begin{itemize}
    \item По освещенности --- рассматриваются видео снятые в светлое время суток, в темное, а также с переходами из темного в светлое (из тоннеля) и из светлого в темное (в тоннель);
    \item По погодным условиям --- рассматриваются видео с ясной, пасмурной, туманной и заснеженной погодой;
    \item По трафику --- рассматриваются видео с не плотным, средним и плотным трафиками;
    \item По окружению --- рассматриваются видео в гороской среде, на трассе и извилистой дороге (серпантине).
\end{itemize}

Тестирование и сравнение параметров всех возможных вариантов для каждой из систем МОТ является довольно тредоемкой и масштабной задачей, так как в результате получим более тысячи видео.

Источниками видеоматериалов являются наборы данных, предназначенные для автономного вождения, такие как MOTChallenge, BDD100K, KITTI. Метрики, по которым будет проводится сравнение перечислены в части%~\ref{subsec:metrics}. 
Таблица с полученными данными приведена в приложении%~\ref{tab::mot_video_results}.

Далее приведены графики сравнения, которые построены на основе полученных результатов. Рассмотрим только случаи, которые представляют наибольший интерес, а именно плотный трафик, темное и изменяющееся время суток, все типы погоды и городская среда.

выбранные условия, на взгяд автора, являются наиболее важными в рамках роботизированного ТС. 
